% Baccalauréat malgache, série D, 2003 : énoncé.
\TCDBacAttribution{D}{2003}
\providecommand{\Courbe}{(\mathscr{C})}
\providecommand{\D}{(\mathscr{D})}
\providecommand{\Plan}{(\mathscr{P})}
\providecommand{\Ouv}{(O ; \vec{u} , \vec{v})}
%\providecommand{\Oij}{(O ; \vec{\imath} , \vec{\jmath})}
\providecommand{\cm}{~\mathrm{cm}}
\providecommand{\dx}{~\mathrm{d}x}
\providecommand{\Arg}[1]{\mathrm{Arg}\left( #1 \right)}
\providecommand{\Angl}[2]{\left( \overrightarrow{#1} ; \overrightarrow{#2} \right)}
%%%%%%%%%%%%%%%%%%%%%%%%%%%%%%%%%%%%%%%%%%%%%%%%%%%%%%%%%%%%%%%%%%%%%%%%%%%%%%%%%%%%%%%%%%%%%%%
%%%%%%%%%%%%%%%%%%%                   %%%%%%%%%%%%%%%%%%%                   %%%%%%%%%%%%%%%%%%%
%%%%%%%%%%%%%%%%%%%%%%%%%%%%%%%%%%%%%%%%%%%%%%%%%%%%%%%%%%%%%%%%%%%%%%%%%%%%%%%%%%%%%%%%%%%%%%%
\noindent \underline{\textsc{Exercice I}} (5 points) \newline
      On considère deux dés cubiques $D$ et $D'$. Le dé $D$, numéroté de $1$ à $6$, est \emph{pipé} de telle sorte que les
      probabilités $p_i$ d'apparition de la face numérotée $i$ vérifient les conditions suivantes :
      \begin{itemize}
            \item $p_1 = p_2 = p_3$,
            \item $p_4 = 3 p_1$,
            \item $p_6 = 2 p_5 = 4 p_4$.
      \end{itemize}
      Le d\'e $D'$ est \emph{non pip\'e}, num\'erot\'e : 1, 1, 1, 2, 2, 4.
      On note $p_1'$, $p_2'$ et $p_4'$ les probabilit\'es d'apparitions respectives des faces num\'erot\'ees 1, 2 et 4.
      \begin{enumerate}
            \item Montrer que $p_1 = \frac{1}{24}$. En d\'eduire $p_i$, $ 2 \leq i \leq 6$.
            \item On lance quatre fois de suite le d\'e $D$ d'une fa\c{c}on ind\'ependante.
                  Soit $X$ la variable al\'eatoire \'egale au nombre d'apparitions de la face num\'erot\'ee 5.
                  \begin{enumerate}
                        \item Donner la loi de probabilit\'e de $X$.
                        \item Calculer l'esp\'erance et la variance de $X$.
                  \end{enumerate}
            \item On lance simultan\'ement les d\'es $D$ et $D'$.
                  Calculer les probabilit\'es des \'ev\`enements suivants : 
                  \begin{itemize}
                        \item $A$ : << La somme des num\'eros obtenus est \'egale \`a $5$ >>.
                        \item $B$ : << La somme des num\'eros obtenus est inf\'erieure ou \'egale \`a $7$ >>.
                        \item $C$ : << Les num\'eros obtenus sont identiques >>.
                  \end{itemize}
      \end{enumerate}
%%%%%%%%%%%%%%%%%%%%%%%%%%%%%%%%%%%%%%%%%%%%%%%%%%%%%%%%%%%%%%%%%%%%%%%%%%%%%%%%%%%%%%%%%%%%%%%
%%%%%%%%%%%%%%%%%%%                   %%%%%%%%%%%%%%%%%%%                   %%%%%%%%%%%%%%%%%%%
%%%%%%%%%%%%%%%%%%%%%%%%%%%%%%%%%%%%%%%%%%%%%%%%%%%%%%%%%%%%%%%%%%%%%%%%%%%%%%%%%%%%%%%%%%%%%%%
\noindent \underline{\textsc{Exercice II}} (5 points) \newline
      Soit le polyn\^ome complexe $P$ d\'efini par : $P(z) = z^2 - (1 + i) \; z + 2 + 2 i$.
      \begin{enumerate}
            \item R\'esoudre dans $\mathbb{C}$ l'\'equation $P(z) = 0$.
            \item Dans le plan complexe $\Plan$ muni d'un rep\`ere orthonorm\'e direct $\Ouv$, on consid\`ere les points $M$
                  d'affixe $z = x + iy$ et $M'$ d'affixe $z' = x' + i y'$, $x$, $y$, $x'$, $y'$ sont des r\'eels.
                  Soit $S$ la similitude plane directe qui, \`a tout point $M$, associe le point $M'$ tel que :
                  \[
                        \begin{multiform}
                              x' = - x - y + 3 \\
                              y' = x - y + 1 \;.
                        \end{multiform}
                  \]
                  \begin{enumerate}
                        \item Ecrire l'expression complexe de $S$.
                        \item Donner ses \'el\'ements caract\'eristiques.
                        \item Donner l'image par $S$ des points $A(0,2)$ et $I(1,2)$.
                  \end{enumerate}
            \item Soit $R$ la rotation de centre $A$ et d'angle $\frac{\pi}{2}$.
                  Donner la nature et les \'el\'ements g\'eom\'etriques de $S' = R \circ S$ o\`u $\circ$ est la 
                  composition des applications.
            \item Quel est l'ensemble $(E)$ des points $M$ d'affixe $z$ v\'erifiant $| (-1+i) \; z + 3 + i | = \sqrt{2}$\,?
      \end{enumerate}
%%%%%%%%%%%%%%%%%%%%%%%%%%%%%%%%%%%%%%%%%%%%%%%%%%%%%%%%%%%%%%%%%%%%%%%%%%%%%%%%%%%%%%%%%%%%%%%
%%%%%%%%%%%%%%%%%%%                   %%%%%%%%%%%%%%%%%%%                   %%%%%%%%%%%%%%%%%%%
%%%%%%%%%%%%%%%%%%%%%%%%%%%%%%%%%%%%%%%%%%%%%%%%%%%%%%%%%%%%%%%%%%%%%%%%%%%%%%%%%%%%%%%%%%%%%%%
\noindent \underline{\textsc{Probleme}} (10 points) \newline
      Soit $f$ la fonction telle que $f(x) = 1 - \frac{1}{2} x - \frac{\ln x}{x}$.
      On note $\Courbe$ la courbe repr\'esentative de $f$ dans un rep\`ere orthonorm\'e \Oij{} d'unit\'e $1\cm$.
      \begin{enumerate}
            \item Calculer $f'$ et $f''$, d\'eriv\'ees premi\`ere et seconde de $f$.
            \item \begin{enumerate}
                        \item Etudier la variation de $f'$.
                        \item En d\'eduire le signe de $f'(x)$ suivant les valeurs de $x$.
                  \end{enumerate}
            \item Etudier la variation de $f$.
            \item Montrer que la droite $\D$ d'\'equation $y = - \frac{1}{2} x + 1$ est asymptote oblique \`a $\Courbe$
                  et \'etudier la position de $\Courbe$ par rapport \`a $\D$ sur $]0,+\infty[$.
            \item Montrer qu'il existe un point unique $A$ de $\Courbe$ o\`u la tangente $(T)$ \`a $\Courbe$ est
                  parall\`ele \`a $\D$.
            \item Tracer $\D$, $(T)$ et $\Courbe$.
            \item Calculer en $\cm^2$, l'aire g\'eom\'etrique du domaine plan limit\'e par la courbe $\Courbe$, 
                  la droite $\D$ et les droites d'\'equations respectives $x = 1$ et $x = e$.
            \item On consid\`ere la fonction $g$ d\'efinie sur $[3,+\infty[$ par : $g(x) = -\frac{1}{2} x + 1 - f(x)$.
                  \begin{enumerate}
                        \item Montrer que $g$ est d\'ecroissante.
                        \item En d\'eduire que si $n \leq x \leq n+1$ alors $0 < g(x) \leq g(n)$ o\`u $n$ est un entier naturel
                              sup\'erieur ou \'egal \`a $3$.
                        \item On d\'efinit la suite $(U_n)_{n\geq 3}$ par $U_n = \int_{n}^{n+1} g(x) \dx$.
                              Donner un encadrement de $U_n$, puis en d\'eduire $\lim U_n$.
                  \end{enumerate}
            %
            \begin{center}
                  On donne : $\expo{-3} \simeq 0,05$ et $\expo{-1} \simeq 0,37$.
            \end{center}
      \end{enumerate}








